\chapter{Sisteme diferențiale}

\section{Calcul matriceal direct}

Ca în cazul sistemelor de ecuații liniare, putem scrie un sistem în formă matriceală $X' = A \cdot X$,
unde $A$ este matricea coeficienților.

Distingem mai multe cazuri, în funcție de proprietățile matricei $A$.

\textbf{Dacă $A$ se diagonalizează}, fie $\{\lambda_i\}$ valorile proprii ale sale,
iar $\{u_i\}$ vectorii proprii asociați. Atunci funcțiile:
\[
  \{ x_i(t) = e^{\lambda_i t} \cdot u_i \}_i
\]
dau o bază în spațiul vectorial al soluțiilor. Soluția generală poate fi scrisă,
atunci, ca o combinație liniară a acestora:
\[
  x(t) = \sum_i c_i x_i(t).
\]

Se obișnuiește ca soluția să fie prezentată sub formă matriceală, alcătuind
\textbf{matricea fundamentală} a sistemului: $X(t) = {}^t (e^{\lambda_i t} u_i)$.

\vspace{0.5cm}

\textbf{Exemplu:}
\[
  \begin{cases}
    x' &= x - y \\
    y' &= -4x + y.
  \end{cases}
\]

Matricea sistemului este $A = \begin{pmatrix} 1 & -1 \\ -4 & 1 \end{pmatrix}$.

Obținem $\lambda_1 = 1$ și $\lambda_2 = 3$. Cum acestea au multiplicitatea algebrică
egală cu 1, rezultă că matricea se diagonalizează. Se obțin subspații proprii de dimensiune
1, bazele fiind, de exemplu:
\[
  X_1 = \begin{pmatrix} 1 \\ 2 \end{pmatrix}, \quad
  X_2 = \begin{pmatrix} 1 \\ -2 \end{pmatrix}.
\]

Atunci soluția generală se scrie:
\[
  X(t) = \begin{pmatrix} x(t) \\ y(t) \end{pmatrix} = %
  C_1 e^{\lambda_1 t} \begin{pmatrix} 1 \\ 2 \end{pmatrix} + %
  C_2 e^{\lambda_2 t} \begin{pmatrix} 1 \\ -2 \end{pmatrix}.
\]

Atunci matricea fundamentală se scrie:
\[
  X(t) = \begin{pmatrix}
    C_1 e^{-t} + C_2 e^{3t} \\
    2C_1 e^{-t} - 2C_2 e^{3t}
  \end{pmatrix}.
\]

\vspace{0.5cm}

\textbf{Dacă $A$ nu se diagonalizează}, fie $\lambda_1, \dots, \lambda_k$ valorile proprii,
de multiplicități algebrice $n_1, \dots, n_k$. Atunci se caută soluții de forma:
\[
  x_j(t) = \sum_{i = 1}^k P_{ij}(t) e^{\lambda_i t}, \quad \mathrm{grad}P_{ij} \leq n_i - 1.
\]

Înlocuim în ecuația inițială și găsim $P_{ij}$. Această metodă poartă numele de
\textbf{metoda coeficienților nedeterminați}.

\vspace{0.5cm}

\textbf{Exemplu:}
\[
  \begin{cases}
    x_1' &= x_2 \\
    x_2' &= -x_1 + 2x_2
  \end{cases}
\]

\emph{Soluție:} Matricea sistemului este $A = \begin{pmatrix} 0 & 1 \\ -1 & 2 \end{pmatrix}$,
care are valoarea proprie $\lambda = 1$, de multiplicitate algebrică 2. Se poate verifica faptul
că matricea nu este diagonalizabilă și atunci căutăm soluții de forma:
\[
  \begin{cases}
    x_1(t) &= (A_1 + B_1 t) e^t \\
    x_2(t) &= (A_2 + B_2 t) e^t
  \end{cases}
\]

Înlocuim în sistem și obținem un sistem liniar cu necunoscutele $A_1, A_2, B_1, B_2$,
care va fi dublu nedeterminat. Așadar, soluția se prezintă cu 2 parametri în forma:
\[
  (A_1, A_2, B_1, B_2) \in \{ (\alpha, \alpha + \beta, \beta, \beta) \mid \alpha, \beta \in \RR \}.
\]

Rezultă că matricea fundamentală este:
\[
  X(t) = \begin{pmatrix}
    (\alpha + \beta t) e^t \\
    (\alpha + \beta + \beta t) e^t
  \end{pmatrix}
\]

\vspace{0.5cm}

\textbf{Dacă matricea $A$ nu se diagonalizează, iar $\lambda_i \in \CC$}, lucrăm cu numere complexe.

\textbf{Exemplu:}
\[
  \begin{cases}
    x'_1 &= 2x_1 + x_2 \\
    x'_2 &= -8x_1 - 2x_2
  \end{cases}
\]
știind că soluția conține la $t = 0$ punctul $M_0(1, 0)$.

\emph{Soluție:} Matricea sistemului este $A = \begin{pmatrix} 2 & 1 \\ -8 & 2 \end{pmatrix}$.

Valorile sale proprii sînt $\lambda_{1,2} = \pm 2i$. Atunci subspațiul invariant va fi:
\[
  V_{\lambda_1} = \mathrm{Sp} \{(1, -2 + 2i )\}.
\]

Obținem o soluție generală în forma:
\[
  x_1(t) = e^{2it} \begin{pmatrix} 1 \\ -2 + 2i \end{pmatrix} = %
  (\cos 2t + i \sin 2t) \begin{pmatrix} 1 \\ -2 + 2i \end{pmatrix}.
\]

Din proprietățile algebrice ale numerelor complexe, avem că $x_t(t) = \overline{x_1(t)}$ și atunci
soluția generală va fi $x(t) = c_1 x_1(t) + c_2 x_2(t)$.

Folosind condițiile inițiale, avem că $x_1(0) = 1$ și $x_2(0) = 0$, deci $c_1 = c_2 = 1$.

%%%%%%%%%%%%%%%%%%%%%%%%%%%%%%%%%%%%%%%%%%%%%%%%%%%%%%%%%%%%%%%%%%%%%%

\section{Folosind forma Jordan}
\label{sec:jordan}

\textbf{Exemplu:} Fie matricea:
\[
  A = \begin{pmatrix}
    1 & -1 & 1 \\
    1 & 1 & -1 \\
    0 & -1 & 2 
  \end{pmatrix}
\]
\begin{enumerate}[(a)]
\item Să se calculeze $e^{At}$, pentru $t \in \RR$;
\item Să se determine soluția generală a sistemului $X' = AX$, apoi soluția particulară
  care verifică condiția $x(0) = 1, y(0) = 2, z(0) = 0$.
\end{enumerate}

\emph{Soluție:}

(a) Matricea $A$ are valorile proprii $\lambda_1 = \lambda_2 = 1, \lambda_3 = 2$.
Vectorii proprii (baze în subspațiile invariante) sînt:
\[
  v_1 = (1, 1, 1), \quad v_2 = (1, 0, 1).
\]

Forma canonică Jordan a matricei se obține a fi:
\[
  J = \begin{pmatrix}
    1 & 0 & 0 \\
    1 & 1 & 0 \\
    0 & 0 & 2
  \end{pmatrix},
\]
folosind matricea de trecere
\[
  T = \begin{pmatrix}
    2 & 1 & 1 \\ 0 & 1 & 0 \\ 1 & 1 & 1
  \end{pmatrix}.
\]

Deci $A = TJT^{-1}$.

Știm din teorie că $e^{At} = T \cdot e^{Jt} \cdot T^{-1}$.

Folosim acum dezvoltarea în serie Taylor a funcției exponențiale:
\[
  e^x = \sum_{n \geq 0} \frac{x}{n!},
\]
adevărată și pentru matrice. În plus, dacă matricea este diagonală,
exponențiala sa va conține ex\-po\-nen\-ți\-a\-la pe diagonală, adică, dacă
$M = \mathrm{diag}(x_1, x_2, \dots, x_n)$, atunci
$e^M = \mathrm{diag}(e^{x_1},\dots, e^{x_n})$.

Pentru a folosi această proprietate, ne legăm de descompunerea în blocuri
Jordan a lui $A$. Fie $J_1$ blocul Jordan de dimensiune 2, corespunzător
valorii proprii $\lambda_1 = \lambda_2 = 1$ și $J_2$ blocul Jordan de dimensiune 1
corespunzător valorii proprii $\lambda_3 = 2$. Avem descompunerea pe blocuri și
pentru exponențială:
\[
  e^{Jt} = \begin{pmatrix}
    e^{J_1 t} & 0 \\
    0 & e^{J_2 t}
  \end{pmatrix}
\]

Cum $e^{J_2 t} = e^{2t}$, calculăm celălalt bloc:
\begin{align*}
  e^{J_1 t} &= e^{\begin{pmatrix} t & 0 \\ t & t \end{pmatrix}} \\
            &= e^{\begin{pmatrix} t & 0 \\ 0 & t \end{pmatrix}
                                               \begin{pmatrix} 0 & 0 \\ t & 0 \end{pmatrix}} \\
            &= \begin{pmatrix} e^t & 0 \\ 0 & e^t \end{pmatrix} \cdot
                                              \Bigg[ \begin{pmatrix} 1 & 0 \\ 0 & 1 \end{pmatrix} +
                                                                                  \frac{1}{1!} \begin{pmatrix} 0 & 0 \\ t & 0 \end{pmatrix} \Bigg] \\
            &= \begin{pmatrix}
              e^t & 0 \\
              te^t & e^t 
            \end{pmatrix}
\end{align*}

Am folosit faptul că $\begin{pmatrix} 0 & 0 \\ t & 0 \end{pmatrix}^2 = 0$, așadar
din seria Taylor va fi suficient să păstrăm primii 2 termeni. Calculăm și $T^{-1}$,
obținînd:
\[
  T^{-1} = \begin{pmatrix}
    1 & 0 & -1 \\
    0 & 1 & 0 \\
    -1 & -1 & 2
  \end{pmatrix}
\]
găsim, în fine:
\[
  e^{At} = T \cdot e^{Jt} \cdot T^{-1} =
  \begin{pmatrix}
    e^t(t + 2) - e^{2t} & e^t - e^{2t} & -e^t(t + 2) + 2e^{2t} \\
    te^t & e^t & -te^t \\
    e^t(t + 1) - e^{2t} & e^t - e^{2t} & -e^t(t + 1) + 2e^{2t}
  \end{pmatrix}
\]

(b) Conform teoriei, soluția generală a sistemului este:
\[
  X(t) = e^{At} \cdot \begin{pmatrix}
    C_1 \\ C_2 \\ C_3
  \end{pmatrix},
\]
constantele $C_i$ determinîndu-se din condițiile inițiale. În acest caz:
\[
  X(t) = e^{At} X(0) = e^{At}
  \begin{pmatrix}
    1 \\ 2 \\ 0 
  \end{pmatrix}
\]

\vspace{0.5cm}

\textbf{Exemplu 2}: Să se determine soluția sistemului diferențial liniar neomogen:
\[
  \begin{cases}
    x' &= 2x - y \\
    y' &= -x + 2z + e^t z \\
    z' &= -y + 2,
  \end{cases}
\]
știind că satisface condițiile inițiale $X(0) = \begin{pmatrix} 0 \\ 1 \\ 0 \end{pmatrix}$.

\emph{Soluție:} Fie matricele:
\[
  X = \begin{pmatrix}
    x \\ y \\ z 
  \end{pmatrix}, \quad
  A = \begin{pmatrix}
    2 & -1 & 0 \\
    -1 & 0 & 2 \\
    0 & -1 & 2 
  \end{pmatrix}, \quad
  B(t) = \begin{pmatrix}
    0 \\ e^t \\ 0
  \end{pmatrix}.
\]

Atunci sistemul se scrie matriceal $X' = AX + B(t)$. Valorile proprii ale matricei $A$
sînt $\lambda_1 = \lambda_2 = 1$ și $\lambda_3 = 2$. Avem vectorul propriu $(1, 1, 1)$
bază în $V_{\lambda_{1,2}}$ și vectorul $(2, 0, 1)$ bază în $V_{\lambda_3}$.

Aducem matricea $A$ la forma canonică Jordan, folosind matricea de trecere:
\[
  T = \begin{pmatrix}
    0 & 1 & 2 \\
    -1 & 1 & 0 \\
    0 & 1 & 1
  \end{pmatrix} \Rightarrow
  T^{-1} = \begin{pmatrix}
    -1 & -1 & 2 \\
    -1 & 0 & 2 \\
    1 & 0 & -1
  \end{pmatrix}
\]

Forma canonică Jordan se obține a fi:
\[
  J = \begin{pmatrix}
    1 & 0 & 0 \\
    1 & 1 & 0 \\
    0 & 0 & 2
  \end{pmatrix}.
\]

Putem descompune matricea Jordan în două blocuri, unul de dimensiune 2 și unul de
dimensiune 1, ceea ce ne ajută să calculăm:
\[
  e^{At} = \begin{pmatrix}
    2e^{2t} - (t + 1)e^t & -te^t & 2e^t(t + 1) - 2 e^{2t} \\
    -te^t & (1 - t)e^t & 2te^t \\
    2e^{2t} - (t + 1)e^t & -te^t & 2e^t(t + 1) -2e^{2t}
  \end{pmatrix}
\]

Conform teoriei, obținem soluția problemei Cauchy (ținînd cont de valorile inițiale) sub forma:
\[
  X(t) = e^{At} X(0) + \int_0^t e^{A(t-s)} B(s) ds.
\]

Integrala se calculează integrînd fiecare din elementele matricei și obținem, în final:
\[
  X(t) = e^t \begin{pmatrix}
    -t - \dfrac{t^2}{2} \\
    \\
    1 - \dfrac{t^2}{2} \\
    \\
    -t - \dfrac{t^2}{2}
  \end{pmatrix}
\]

%%% Local Variables:
%%% mode: latex
%%% TeX-master: "../ag-20-21"
%%% End:
